\documentclass[paper=a4,twoside=false,BCOR=0mm,DIV=calc,fontsize=12pt]{scrartcl}

\usepackage[automark,headsepline]{scrpage2}
% \usepackage{xunicode,fontspec,xltxtra}
\usepackage[english,german]{babel}

\usepackage[T1]{fontenc}
\usepackage[utf8x]{inputenc}

\usepackage{graphicx}
\usepackage{lastpage}
\usepackage{listings}

%Schöne Tabellen
\usepackage{booktabs}

%\usepackage{longtable}

\usepackage{multirow}

% Kein einrücken und Abstand zwischen Absätzen.
\usepackage{parskip}
\setlength{\parindent}{0cm}

% Hyperlinks und url's
\usepackage[hidelinks]{hyperref}



% \usepackage[debug]{libertine}
% \setromanfont[Mapping=tex-text]{Linux Libertine O}
% \setsansfont[Mapping=tex-text]{Linux Biolinum O}
% \setmonofont[Mapping=tex-text]{Liberation Mono}

\pagestyle{scrheadings}

\clearscrheadfoot
\ihead{\headmark}
\ohead{Seite\pagemark\ von \pageref{LastPage}}
\ifoot{Business frameworks}
\cfoot{{\includegraphics[width=1.5cm]{./img/CC3.png}}}
\ofoot{\today}

\renewcommand*{\pnumfont}{
	\normalfont\rmfamily\slshape
}

\KOMAoptions{draft=false}
\KOMAoptions{DIV=last}

\begin{document}

% --- Titelseite --- %
\begin{titlepage}
	\enlargethispage{3cm}
	\begin{raggedright}
	\begin{picture}(0,0)
		\put(-30,14){\includegraphics[width=7cm]{./img/fhnw-technik-head}}
	\end{picture}

	\vspace*{6cm}
	{\Huge\bfseries\sf
		MAS-IT\\[1.7ex]
	}
	{\Large\bfseries\sf
		Business frameworks\\[2.2ex]
	}
	{\large\bfseries\sf
		Masterthesis von\\[1.5ex]
		Etienne Rebetez\\[1.5ex]
	}
	\vspace*{1.5cm}
	{\large\bfseries\sf
		FHNW\\[1.5ex]
		Hochschule für Technik\\[1.5ex]
		Studiengang MAS-IT\\[1.5ex]
		Betreuender Dozent:\\[1.5ex]
		Dr. sc. techn. Ronald Tanner\\[1.5ex]
	}
	\vspace*{2cm}
	{\large\bfseries\sf
		Bern, \today\\
	}
	\end{raggedright}
\end{titlepage}

\newpage
% --- Zusammenfassung --- %
\section*{Summary}
Text

% --- Vorwort --- %
\section*{Vorwort}
Text
% --- Inhaltsverzeichnis --- %
\newpage
	\tableofcontents

% --- Einleitung --- %
\newpage
\section{Einleitung}
Laborprozesse. Was versteht man eigentlich darunter? 
Der stark abstrahierte Fall ist der folgende. Ein Physische Probe kommt in ein Analyselabor, bestimmte physikalische Parameter werden 
von der Probe erhoben und diese Resultate werden dem Auftraggeber bekanntgegeben.

In diesem Bericht wird ein Labor als Qualitätsicherungslabor betrachtet. Das heisst, die Prozesse und Methoden sind vorgeben und können bzw. dürfen nicht 
einfach abgeändert werden.

Traditionell wird die Probenverwaltung von einem LIMS (Labor Informations und Management System) geregelt. Traditionellerweise ist es auch so, dass
sobald die Probe welche analysiert werden muss im Labor steht, viele Schritte Manuell geschehen. Dies ist aufwendig und Fehleranfällig. 
Es darf natürlich nicht vergessen werden, dass unter Manuell auch Handarbeiten wie einwägen, Lösungen erstellen oder Analysegeräte befüllen fallen.
Für diese Aufgaben wird es noch lange Menschen geben welche diese Arbeit ausführen müssen.

Will man nun diese Laborprozesse Automatisieren, geht es vor allem um die Datenübertragung. 
Was muss gemessen werden, welche Methode soll verwendet werden,
welches qualifizierte Gerät darf verwendet werden, 
wie werden die Daten verrechnet und schlussendlich wie werden die Resultate wieder an das nächst höhere System weitergegeben.
Gleichzeitig solle das System den Laborarbeiter unterstützen und die durchgeführten Schritte sollten in alle nachvollziehbar sein.

Im Labor gibt es nun verschiedenste Methoden welche unterschiedliche Ursprünge haben. So ist eine Chemische Titration ein vollkommen anderes Konzept 
als z.B eine Biochemische Elysa Analyse.
Dies ergibt ein äusserst Heterogene Landschaft an Analysegeräten mit und ohne Software. Resultate oder Rohdaten liegen in unterschiedlichsten Formaten und Layouts vor.

<Geräte Software>



Menschen sind ausserordentlich gut darin Informationen 







\section{Konzept}


\section{BPMN}


\section{BPEL}



\section{Engines}

\subsection{activiti}


\subsection{activiti}






\addcontentsline{toc}{section}{Literaturverzeichnis}
\begin{thebibliography}{99}

\bibitem{activiti} \url{http://www.activiti.org/}




\end{thebibliography}

\end{document}

